\documentclass[../main.tex]{subfiles}

\begin{document}

\chapter{不等式}\label{chapter:appB-不等式}

\section{凸函数}\label{section:appB-凸函数}

\begin{definition}\label{definition:appB-凸函数}
函数 $f:\mathbb{R}^n\to\mathbb{R}$ 称为凸的，只要
\begin{align}
f(\tau x+(1-\tau)y)\leqslant\tau f(x)+(1-\tau)f(y)
\label{align:appB-凸性}
\end{align}
对所有 $x,y\in\mathbb{R}^n$ 和每个 $0\leqslant\tau\leqslant 1$。
\end{definition}

\begin{theorem}[支撑超平面。]\label{theorem:appB-支撑超平面}设 $f:\mathbb{R}^n\to\mathbb{R}$ 凸。则对每个 $x\in\mathbb{R}^n$ 存在 $r\in\mathbb{R}^n$ 使得不等式
\begin{align}
f(y)\geqslant f(x)+r\cdot(y-x)
\label{align:appB-支撑}
\end{align}
对所有 $y\in\mathbb{R}^n$ 成立。
\end{theorem}

映射 $y\mapsto f(x)+r\cdot(y-x)$ 确定了 $f$ 在 $x$ 处的支撑超平面。不等式\eqref{align:appB-支撑} 说 $f$ 的图像位于每个支撑超平面之上。若 $f$ 在 $x$ 处可微，$r=Df(x)$。

若 $f$ 是 $C^2$ 的，则 $f$ 凸当且仅当 $D^2f\geqslant 0$。$C^2$ 函数 $f$ 一致凸，若 $D^2f\geqslant\theta I$ 对某常数 $\theta>0$：这意味着
\begin{align*}
\sum_{i,j=1}^n f_{x_ix_j}(x)\xi_i\xi_j\geqslant\theta|\xi|^2 \quad (x,\xi\in\mathbb{R}^n).
\end{align*}

\begin{theorem}[Jensen 不等式。]\label{theorem:appB-Jensen}设 $f:\mathbb{R}^m\to\mathbb{R}$ 凸，$U\subset\mathbb{R}^n$ 开、有界。设 $\mathbf{u}:U\to\mathbb{R}^m$ 可积。则
\begin{align}
f\left(\fint_U\mathbf{u}\,\mathrm{d}x\right)\leqslant\fint_U f(\mathbf{u})\,\mathrm{d}x.
\label{align:appB-Jensen}
\end{align}
\end{theorem}

回忆 \S A.3 中记号 $\fint_U\mathbf{u}\,\mathrm{d}x=\dfrac{1}{|U|}\int_U\mathbf{u}\,\mathrm{d}x=\mathbf{u}$ 在 $U$ 上的平均。

\begin{proof}
由于 $f$ 凸，对每个 $p\in\mathbb{R}^m$ 存在 $r\in\mathbb{R}^m$ 使得
\begin{align*}
f(q)\geqslant f(p)+r\cdot(q-p)
\end{align*}
对所有 $q\in\mathbb{R}^m$。

令 $p=\fint_U\mathbf{u}\,\mathrm{d}y$，$q=\mathbf{u}(x)$：
\begin{align*}
f(\mathbf{u}(x))\geqslant f\left(\fint_U\mathbf{u}\,\mathrm{d}y\right)+r\cdot\left(\mathbf{u}(x)-\fint_U\mathbf{u}\,\mathrm{d}y\right).
\end{align*}
对 $x$ 在 $U$ 上积分。
\end{proof}

我们在 \S3.3.2 和 \S9.6.1 中更充分地讨论凸函数。

\section{有用的不等式}\label{section:appB-有用不等式}

以下是一组初等但基本的不等式。这些估计贯穿全书不断使用，应当记住。

\textbf{a. Cauchy 不等式。}
\begin{align}
ab\leqslant\frac{a^2}{2}+\frac{b^2}{2} \quad (a,b\in\mathbb{R}).
\label{align:appB-Cauchy}
\end{align}
\begin{proof}
$0\leqslant(a-b)^2=a^2-2ab+b^2$。
\end{proof}

\begin{theorem}[带 $\varepsilon$ 的 Cauchy 不等式]\label{theorem:带参的Cauchy不等式}
\begin{align}
ab\leqslant\varepsilon a^2+\frac{b^2}{4\varepsilon} \quad (a,b>0,\;\varepsilon>0).
\label{align:appB-CauchyEps}
\end{align}
\end{theorem}
\begin{proof}
写 $ab=((2\varepsilon)^{1/2}a)\left(\dfrac{b}{(2\varepsilon)^{1/2}}\right)$ 并应用 Cauchy 不等式。
\end{proof}

\textbf{c. Young 不等式。} 设 $1<p,q<\infty$，$\dfrac{1}{p}+\dfrac{1}{q}=1$。则
\begin{align}
ab\leqslant\frac{a^p}{p}+\frac{b^q}{q} \quad (a,b>0).
\label{align:appB-Young}
\end{align}
\begin{proof}
映射 $x\mapsto e^x$ 凸，因此
\begin{align*}
ab=e^{\log a+\log b}=e^{\frac{1}{p}\log a^p+\frac{1}{q}\log b^q}\leqslant\frac{1}{p}e^{\log a^p}+\frac{1}{q}e^{\log b^q}=\frac{a^p}{p}+\frac{b^q}{q}.
\end{align*}
\end{proof}

\textbf{d. 带 $\varepsilon$ 的 Young 不等式。}
\begin{align}
ab\leqslant\varepsilon a^p+C(\varepsilon)b^q \quad (a,b>0,\;\varepsilon>0)
\label{align:appB-YoungEps}
\end{align}
其中 $C(\varepsilon)=(\varepsilon p)^{-q/p}q^{-1}$。
\begin{proof}
写 $ab=((\varepsilon p)^{1/p}a)\left(\dfrac{b}{(\varepsilon p)^{1/p}}\right)$ 并应用 Young 不等式。
\end{proof}

\textbf{e. H\"older 不等式。} 设 $1\leqslant p,q\leqslant\infty$，$\dfrac{1}{p}+\dfrac{1}{q}=1$。则若 $u\in L^p(U)$，$v\in L^q(U)$，有
\begin{align}
\int_U|uv|\,\mathrm{d}x\leqslant\|u\|_{L^p(U)}\|v\|_{L^q(U)}.
\label{align:appB-Holder}
\end{align}
\begin{proof}
由齐次性，可假设 $\|u\|_{L^p}=\|v\|_{L^q}=1$。则对 $1<p,q<\infty$，Young 不等式蕴含
\begin{align*}
\int_U|uv|\,\mathrm{d}x\leqslant\frac{1}{p}\int_U|u|^p\,\mathrm{d}x+\frac{1}{q}\int_U|v|^q\,\mathrm{d}x=1=\|u\|_{L^p}\|v\|_{L^q}.
\end{align*}
\end{proof}

\textbf{f. Minkowski 不等式。} 设 $1\leqslant p\leqslant\infty$，$u,v\in L^p(U)$。则
\begin{align}
\|u+v\|_{L^p(U)}\leqslant\|u\|_{L^p(U)}+\|v\|_{L^p(U)}.
\label{align:appB-Minkowski}
\end{align}
\begin{proof}
\begin{align*}
\|u+v\|_{L^p(U)}^p &=\int_U|u+v|^p\,\mathrm{d}x\\
&\leqslant\int_U|u+v|^{p-1}(|u|+|v|)\,\mathrm{d}x\\
&\leqslant\left(\int_U|u+v|^p\,\mathrm{d}x\right)^{\frac{p-1}{p}}
\left[\left(\int_U|u|^p\,\mathrm{d}x\right)^{1/p}
+\left(\int_U|v|^p\,\mathrm{d}x\right)^{1/p}\right]\\
&=\|u+v\|_{L^p(U)}^{p-1}(\|u\|_{L^p(U)}+\|v\|_{L^p(U)}).
\end{align*}
\end{proof}

\begin{remark}\label{remark:appB-离散版本}
类似证明建立 H\"older 和 Minkowski 不等式的离散版本：
\begin{align}
\left|\sum_{k=1}^n a_kb_k\right|\leqslant\left(\sum_{k=1}^n|a_k|^p\right)^{1/p}\left(\sum_{k=1}^n|b_k|^q\right)^{1/q},
\label{align:appB-离散Holder}
\end{align}
\begin{align*}
\left(\sum_{k=1}^n|a_k+b_k|^p\right)^{1/p}\leqslant\left(\sum_{k=1}^n|a_k|^p\right)^{1/p}+\left(\sum_{k=1}^n|b_k|^p\right)^{1/p},
\end{align*}
其中 $a=(a_1,\cdots,a_n)$，$b=(b_1,\cdots,b_n)\in\mathbb{R}^n$，$1\leqslant p<\infty$，$\dfrac{1}{p}+\dfrac{1}{q}=1$。
\end{remark}

\begin{theorem}[广义H\"older 不等式]\label{theorem:广义Holder不等式}
设 $1\leqslant p_1,\cdots,p_m\leqslant\infty$，$\dfrac{1}{p_1}+\dfrac{1}{p_2}+\cdots+\dfrac{1}{p_m}=1$，假设 $u_k\in L^{p_k}(U)$（$k=1,\cdots,m$）。则
\begin{align}
\int_U|u_1\cdots u_m|\,\mathrm{d}x\leqslant\prod_{k=1}^m\|u_k\|_{L^{p_k}(U)}.
\label{align:appB-广义Holder}
\end{align}
\end{theorem}
\begin{proof}
归纳法，使用 H\"older 不等式。
\end{proof}

\textbf{h. $L^p$ 范数的插值不等式。} 设 $1\leqslant s\leqslant r\leqslant t\leqslant\infty$，
\begin{align*}
\frac{1}{r}=\frac{\theta}{s}+\frac{1-\theta}{t}.
\end{align*}
又设 $u\in L^s(U)\cap L^t(U)$。则 $u\in L^r(U)$，且
\begin{align}
\|u\|_{L^r(U)}\leqslant\|u\|_{L^s(U)}^\theta\|u\|_{L^t(U)}^{1-\theta}.
\label{align:appB-插值}
\end{align}
\begin{proof}
计算
\begin{align*}
\int_U|u|^r\,\mathrm{d}x=\int_U|u|^{\theta r}|u|^{(1-\theta)r}\,\mathrm{d}x\leqslant\left(\int_U|u|^{\frac{\theta r}{\theta r/s}}\,\mathrm{d}x\right)^{\frac{\theta r}{s}}\left(\int_U|u|^{\frac{(1-\theta)r}{(1-\theta)r/t}}\,\mathrm{d}x\right)^{\frac{(1-\theta)r}{t}}.
\end{align*}
我们应用了 H\"older 不等式，其适用因 $\dfrac{\theta r}{s}+\dfrac{(1-\theta)r}{t}=1$。
\end{proof}

\textbf{i. Cauchy--Schwarz 不等式。}
\begin{align}
|x\cdot y|\leqslant|x||y| \quad (x,y\in\mathbb{R}^n).
\label{align:appB-CauchySchwarz}
\end{align}
\begin{proof}
设 $\varepsilon>0$，注意
\begin{align*}
0\leqslant|x\pm\varepsilon y|^2=|x|^2\pm 2\varepsilon x\cdot y+\varepsilon^2|y|^2.
\end{align*}
因此
\begin{align*}
\pm x\cdot y\leqslant\frac{1}{2\varepsilon}|x|^2+\frac{\varepsilon}{2}|y|^2.
\end{align*}
令 $\varepsilon=|x|/|y|$ 最小化右端（设 $y\neq 0$）。
\end{proof}

\begin{remark}\label{remark:appB-矩阵CS}
类似地，若 $A$ 是对称非负定 $n\times n$ 矩阵，
\begin{align}
\left|\sum_{i,j=1}^n a_{ij}x_iy_j\right|\leqslant\left(\sum_{i,j=1}^n a_{ij}x_ix_j\right)^{1/2}\left(\sum_{i,j=1}^n a_{ij}y_iy_j\right)^{1/2} \quad (x,y\in\mathbb{R}^n).
\label{align:appB-矩阵CS}
\end{align}
\end{remark}

\textbf{j. Gronwall 不等式（微分形式）。}
\begin{enumerate}[(i)]
\item 设 $\eta(\cdot)$ 是 $[0,T]$ 上非负绝对连续函数，对 a.e. $t$ 满足微分不等式
\begin{align}
\eta'(t)\leqslant\phi(t)\eta(t)+\psi(t),
\label{align:appB-Gronwall微分}
\end{align}
其中 $\phi(t)$ 和 $\psi(t)$ 是 $[0,T]$ 上非负可积函数。则
\begin{align}
\eta(t)\leqslant e^{\int_0^t\phi(s)\,\mathrm{d}s}\left[\eta(0)+\int_0^t\psi(s)\,\mathrm{d}s\right]
\label{align:appB-Gronwall结果}
\end{align}
对所有 $0\leqslant t\leqslant T$。

\item 特别地，若 $\eta'\leqslant\phi\eta$ 在 $[0,T]$ 上且 $\eta(0)=0$，则 $\eta\equiv 0$ 在 $[0,T]$ 上。
\end{enumerate}

\begin{proof}
由\eqref{align:appB-Gronwall微分} 得
\begin{align*}
\frac{\mathrm{d}}{\mathrm{d}s}\left(\eta(s)e^{-\int_0^s\phi(r)\,\mathrm{d}r}\right)=e^{-\int_0^s\phi(r)\,\mathrm{d}r}(\eta'(s)-\phi(s)\eta(s))\leqslant e^{-\int_0^s\phi(r)\,\mathrm{d}r}\psi(s)
\end{align*}
对 a.e. $0\leqslant s\leqslant T$。因此对每个 $0\leqslant t\leqslant T$，
\begin{align*}
\eta(t)e^{-\int_0^t\phi(r)\,\mathrm{d}r}\leqslant\eta(0)+\int_0^t e^{-\int_0^s\phi(r)\,\mathrm{d}r}\psi(s)\,\mathrm{d}s\leqslant\eta(0)+\int_0^t\psi(s)\,\mathrm{d}s.
\end{align*}
这蕴含不等式\eqref{align:appB-Gronwall结果}。
\end{proof}

\end{document}